\section{Eigenvectors and Eigenvalues}
\label{sec:eigen}

\noindent\textbf{Subject:} Linear Algebra for ML (3Blue1Brown) \\
\textbf{Status:} growing

\subsection*{Summary}
Eigenvectors are special vectors that stay on their own span after a linear transformation. Eigenvalues are the factor by which they get stretched or squished.

\subsection{Eigenvectors}
\begin{defbox}
Any vector whose \textbf{span stays the same} after the linear transformation.
\end{defbox}

\subsection{Eigenvalues}
\begin{defbox}
Just the \textbf{factor by which it's stretched or squished} during the linear transformation.
\end{defbox}

\subsection{Symbolically}

\[
A\vec{v} = \lambda\vec{v} \quad \rightarrow \text{Eigenvalue}
\]

\begin{itemize}[nosep]
  \item $A$ = transformation matrix (matrix-vector multiplication)
  \item $\vec{v}$ = eigenvector
  \item $\lambda$ = eigenvalue (scalar)
\end{itemize}

\begin{notebox}
It says that the \textbf{matrix-vector multiplication} gives the same result as just \textbf{scaling the eigenvector} $\vec{v}$ by some value $\lambda$.
\end{notebox}

\subsection{Seeking Eigenvalue $\lambda$}

\[
A\vec{v} = \lambda\vec{v}
\]
\[
A\vec{v} - \lambda I\vec{v} = 0
\]
\[
(A - \lambda I)\vec{v} = 0
\]
\[
\det(A - \lambda I) = 0
\]

Matrix multiplication by $\lambda$:

\[
\begin{bmatrix} \lambda & 0 & 0 \\ 0 & \lambda & 0 \\ 0 & 0 & \lambda \end{bmatrix}
\Rightarrow
\begin{bmatrix} 1 & 0 & 0 \\ 0 & 1 & 0 \\ 0 & 0 & 1 \end{bmatrix}
\]

We want a \textbf{non-zero solution} for $\vec{v}$:
\begin{itemize}[nosep]
  \item ``Squishification'' $\rightarrow$ $\det(A - \lambda I) = 0$
  \item This matrix will look like: $\begin{bmatrix} 3-\lambda & 1 & 4 \\ 1 & 5-\lambda & 9 \\ 2 & 6 & 5-\lambda \end{bmatrix}$
\end{itemize}

$\rightarrow$ So the $\lambda$ value should be 1, meaning this value keeps the vector on a line.

\paragraph{Steps:}
\begin{enumerate}[nosep]
  \item $A\vec{v} = \lambda\vec{v}$
  \item $A\vec{v} - \lambda I\vec{v} = 0$
  \item $(A - \lambda I)\vec{v} = 0$
  \item $\det(A - \lambda I) = 0$
\end{enumerate}

\paragraph{Example:}
\[
\det\left(\begin{bmatrix} 3-\lambda & 1 \\ 0 & 2-\lambda \end{bmatrix}\right) = (3-\lambda)(2-\lambda) - 0 \cdot 1
\]

This is a \textbf{quadratic polynomial in $\lambda$}.

\[
(3-\lambda)(2-\lambda) = 0
\]
\[
\lambda = 3 \quad \text{or} \quad \lambda = 2
\]

\begin{notebox}
Some vectors don't have an eigenvalue (e.g.\ a 90\textdegree{} rotation).
\end{notebox}

\subsection{Or Complex Numbers}

\[
\det\left(\begin{bmatrix} -\lambda & 1 \\ 1 & -\lambda \end{bmatrix}\right) = (-\lambda)(-\lambda) - 1 = \lambda^2 - 1 = 0
\]
\[
\lambda = i \quad \text{or} \quad \lambda = -i
\]

\subsection{Eigenbasis}
\begin{defbox}
It's the same concept as the \textbf{change of basis}.
\end{defbox}

We use our eigenvectors as the (new) basis --- \textbf{Diagonal}:

\[
\begin{bmatrix} 1 & -1 \\ 0 & 1 \end{bmatrix}^{-1}
\begin{bmatrix} 3 & 1 \\ 0 & 2 \end{bmatrix}
\begin{bmatrix} 1 & -1 \\ 0 & 1 \end{bmatrix}
=
\begin{bmatrix} 3 & 0 \\ 0 & 2 \end{bmatrix}
\]

$\rightarrow$ Change of basis matrix

\subsection{Some Facts to Know}

\paragraph{1) Trace:}
\[
\operatorname{Tr}\left(\begin{bmatrix} a & b \\ c & d \end{bmatrix}\right) = a + d = \lambda_1 + \lambda_2
\]
\[
\frac{1}{2}\operatorname{Tr}\left(\begin{bmatrix} a & b \\ c & d \end{bmatrix}\right) = \frac{a+d}{2} = \frac{\lambda_1 + \lambda_2}{2}
\]
\[
\text{mean}(a,d) = \text{mean}(\lambda_1, \lambda_2)
\]

\paragraph{2) Determinant:}
\[
\det\left(\begin{bmatrix} a & b \\ c & d \end{bmatrix}\right) = ad - bc = \lambda_1 \lambda_2
\]

\paragraph{3) Eigenvalue shortcut:}
\[
\lambda_1, \lambda_2 = m \pm \sqrt{m^2 - p}
\]
where $m$ = mean of the diagonal, $p$ = product (determinant).

\subsection{Finding Eigenvalues --- Example}

We call the mean $(m)$ the trace, and the determinant $(p)$ the product.

\paragraph{Find the distance:}
\[
\begin{bmatrix} 8 & 4 \\ 2 & 6 \end{bmatrix} \quad m = 7, \quad p = 40
\]

\[
40 = (7-d)(7+d)
\]
\[
40 = 7^2 - d^2
\]
\[
d^2 = 7^2 - 40 \Rightarrow d^2 = 9 \Rightarrow d = 3
\]

So:
\[
\lambda_{1,2} = m \pm d = 7 \pm 3 \Rightarrow \lambda_1 = 4, \quad \lambda_2 = 10
\]

\paragraph{General formula:}
\[
d^2 = m^2 - p \quad \Rightarrow \quad \lambda_{1,2} = m \pm \sqrt{m^2 - p}
\]

\paragraph{Example:} $x^2 - 10x + 9 = 5$, so for $x^2 - 10x + 9$:
\begin{itemize}[nosep]
  \item Mean = 5
  \item $(x - \lambda_1)(x - \lambda_2)$
  \item $x^2 - (\lambda_1 + \lambda_2)x + \lambda_1\lambda_2$ $\rightarrow$ $m$ (sum), product
\end{itemize}

\subsection{Find the Eigenvalues of $\begin{bmatrix} 3 & 1 \\ 4 & 1 \end{bmatrix}$}

\[
\det\left(\begin{bmatrix} 3-\lambda & 1 \\ 4 & 1-\lambda \end{bmatrix}\right) = (3-\lambda)(1-\lambda) - 4(1)
\]
\[
= (3 - 3\lambda - \lambda + \lambda^2) - 4
\]
\[
= (3 - 4\lambda + \lambda^2) - 4
\]
\[
= \lambda^2 - 4\lambda - 1 = 0
\]

\[
\lambda_{1,2} = \frac{4 \pm \sqrt{4^2 - 4(1)(-1)}}{2} = \frac{4 \pm \sqrt{20}}{2} = 2 \pm \sqrt{5}
\]

\subsection*{Related Topics}
\begin{itemize}[nosep]
  \item Section~\ref{sec:change-of-basis} --- Change of Basis
  \item Section~\ref{sec:matrices} --- Matrices
\end{itemize}

\bigskip
\noindent\textit{Tags: linear-algebra, eigenvectors, eigenvalues, eigenbasis, math, phase-1}
